\documentclass[11pt]{article}
\usepackage[margin=1in]{geometry}
\usepackage{amsmath,amssymb}
\usepackage{float}
\usepackage{tabularx}
\usepackage{hyperref}

\title{Function and Parameter Reference: \texttt{model\_Bacteremia.py}}
\author{}
\date{}

\begin{document}
\maketitle

\section{Overview}

\texttt{model\_Bacteremia.py} is the base simulation model imported by every
other script in this project. It defines two classes --
\texttt{PharmacokineticModel} and \texttt{ImmuneResponse} -- and one
standalone function, \texttt{dual\_reservoir\_model}, which together
describe drug exposure, host immune clearance, and bacterial population
dynamics for a two-strain, two-compartment (blood and bone reservoir)
infection. A bottom-of-file \texttt{\_\_main\_\_} block runs one
demonstration simulation when the file is executed directly; it is not a
function and is not used by any other script, which instead import the
class and function definitions directly. This document lists every
function, its arguments, what it does, and the parameter values used
throughout the project -- including, where applicable, the coefficient of
variation (CV) used to sample that parameter in the project's Monte Carlo
scripts. No equations are given here; see the project's Results/Methods
write-up for the underlying mathematical relationships.

\section{Class: \texttt{PharmacokineticModel}}

Represents the two-drug dosing regimen (vancomycin, then linezolid) and
converts it into a free-drug concentration profile over time.

\subsection{\texttt{\_\_init\_\_(self)}}

\textbf{Arguments:} none (besides \texttt{self}).

\textbf{What it does:} Sets fixed dosing and pharmacokinetic constants for
both drugs as instance attributes (dose, dosing interval, course duration,
volume of distribution, free/unbound fraction, and clearance). These
values are imported from the centralized \texttt{PK\_PARAMS} mapping in
\texttt{parameters.py} and assigned as instance attributes
(Table~\ref{tab:pk}).

\subsection{\texttt{concentration\_function(self, drug\_type, total\_time\_h, start\_h=0)}}

\textbf{Arguments:}
\begin{itemize}
\item \texttt{drug\_type} (\texttt{str}) -- \texttt{'vancomycin'} selects
the vancomycin constants; any other value selects linezolid.
\item \texttt{total\_time\_h} (\texttt{float}) -- total length, in hours, of
the time grid over which the concentration profile is built. Must cover the
full simulated course.
\item \texttt{start\_h} (\texttt{float}, default \texttt{0}) -- the
simulation time, in hours, at which the first dose of this drug is given.
\end{itemize}

\textbf{What it does:} Builds a repeated-dose, first-order-elimination
concentration-time profile for the selected drug on a dense grid (10 points
per hour), by superposing one decaying exponential per dose given between
\texttt{start\_h} and \texttt{start\_h + duration}. Concentration is forced
to zero for all times before \texttt{start\_h}. Returns a single-argument
callable, \texttt{conc\_func(t)}, that linearly interpolates this profile at
any query time \texttt{t}. This returned function is what gets passed into
\texttt{dual\_reservoir\_model} as \texttt{van\_func} or \texttt{lzd\_func}.

\begin{table}[H]
\centering
\small
\renewcommand{\arraystretch}{1.15}
\caption{Fixed values set by \texttt{PharmacokineticModel.\_\_init\_\_}.}
\vspace{0.5cm}
\begin{tabularx}{\textwidth}{l l l X}
\hline
\textbf{Parameter} & \textbf{Vancomycin} & \textbf{Linezolid} & \textbf{Meaning} \\
\hline
Dose & 1000 mg & 600 mg & amount per administration \\
Interval & 8 h & 12 h & time between doses (q8h / q12h) \\
Duration & 96 h (4 d) & 1008 h (42 d) & total length of the drug course \\
$V_d$ & 50 L & 45 L & volume of distribution \\
$f_u$ & 0.50 & 0.70 & free (unbound, active) fraction \\
Clearance & 5.0 L/h & 6.0 L/h & total body clearance \\
\hline
\end{tabularx}
\label{tab:pk}
\end{table}

\section{Class: \texttt{ImmuneResponse}}

Represents the host's bacterial clearance capacity as a constant rate,
independently scaled per compartment.

\subsection{\texttt{\_\_init\_\_(self, k\_immune=K\_IMMUNE\_DEFAULT, eff\_blood=1.0, eff\_res=0.5)}}

\textbf{Arguments:}
\begin{itemize}
\item \texttt{k\_immune} (\texttt{float}, default
\texttt{K\_IMMUNE\_DEFAULT} $=0.18$ from \texttt{parameters.py}) -- the
baseline host immune clearance rate, in h$^{-1}$.
\item \texttt{eff\_blood} (\texttt{float}, default \texttt{1.0}) -- a
multiplier scaling how effectively the immune response clears bacteria in
the blood compartment.
\item \texttt{eff\_res} (\texttt{float}, default \texttt{0.5}) -- the
equivalent multiplier for the bone reservoir compartment. Set to one-half
of \texttt{eff\_blood} to represent reduced immune access to a deep-seated
focus of infection.
\end{itemize}

\textbf{What it does:} Stores these three values as instance attributes for
later use by \texttt{dual\_reservoir\_model}.

\subsection{\texttt{compute(self, t=None)}}

\textbf{Arguments:} \texttt{t} (optional, default \texttt{None}) -- accepted
but not used.

\textbf{What it does:} Returns \texttt{self.k\_immune} unconditionally. The
unused \texttt{t} argument exists purely so this method can be called with
the same signature as a time-dependent function (matching how
\texttt{van\_func(t)}/\texttt{lzd\_func(t)} are called), even though the
current immune response does not actually vary with time.

\section{Function: \texttt{dual\_reservoir\_model(y, t, params, van\_func, lzd\_func, immune\_model)}}

The right-hand side of the ordinary differential equation (ODE) system,
written in the form consumed by the project's shared
\texttt{integrate\_model} solver.

\textbf{Arguments:}
\begin{itemize}
\item \texttt{y} -- length-4 array/list of the current state,
$[S_b,\ R_b,\ S_{res},\ R_{res}]$: sensitive and resistant \textit{S.
aureus} counts (CFU/mL) in blood ($b$) and the bone reservoir ($res$).
\item \texttt{t} -- the current simulation time (h). Passed through to
\texttt{van\_func}, \texttt{lzd\_func}, and \texttt{immune\_model.compute}
so that time-dependent drug exposure (and, if ever implemented, a
time-dependent immune response) can be evaluated at this instant.
\item \texttt{params} -- a dictionary of growth, drug-effect, capacity, and
exchange-rate values (Table~\ref{tab:params}). Analysis scripts copy the
shared \texttt{BASE\_PARAMS} mapping and override only the parameter under
study.
\item \texttt{van\_func}, \texttt{lzd\_func} -- callables returning free
vancomycin/linezolid concentration at time \texttt{t}, normally produced by
\texttt{PharmacokineticModel.concentration\_function}.
\item \texttt{immune\_model} -- an \texttt{ImmuneResponse} instance.
\end{itemize}

\textbf{What it does:} Computes, from the current state and the drug/immune
inputs, the instantaneous rate of change of each of the four bacterial
subpopulations -- combining logistic growth toward a compartment-specific
carrying capacity, drug-driven killing (vancomycin against sensitive
strains only; linezolid against both strains, at a discounted effect in the
reservoir), immune-driven killing (at a discounted rate in the reservoir),
and one-way first-order seeding of bacteria from the reservoir into blood
(there is no blood-to-reservoir return term). Non-physical negative populations are clamped to zero.
Reservoir linezolid exposure is multiplied by
\texttt{lzd\_res\_fraction}; its reservoir maximum effect is additionally
multiplied by the independent efficacy factor \texttt{lzd\_res\_efficacy}
($\alpha$), so $E_{max,L}^{res} = \alpha\,E_{max,L}$. The factor is not
derived from the growth rates (it replaced the former $\rho_{res,S}/\rho_S$
scaling). This function is the only place $\alpha$ is used; every analysis
script applies it through \texttt{BASE\_PARAMS}, and none varies it. Growth otherwise follows the ordinary
logistic term for active states, allowing a detectable reservoir population
to reseed blood.

The numerical integration applies a hard limit of detection (LOD) of
$10$~CFU/mL: a compartment that falls below the LOD is set to exactly zero.
Thus, sub-LOD fractional populations cannot subsequently regrow, although a
nonzero reservoir can still reseed a cleared blood compartment through the
exchange term.

\textbf{Returns:} a length-4 list, $[\dot S_b,\ \dot R_b,\ \dot S_{res},\
\dot R_{res}]$, the derivatives the shared solver integrates forward in time.

\section{Shared numerical solver}

The model right-hand side does not call SciPy directly. Analysis scripts pass
it to \texttt{solver.integrate\_model}, which uses
\texttt{scipy.integrate.solve\_ivp} with LSODA. The centralized numerical
configuration is relative tolerance $10^{-6}$, absolute tolerance $10^{-8}$
CFU/mL, and maximum internal step 4.0~h. The standard output grid contains
1453 points and includes all treatment boundaries exactly.

The solver implements extinction and reactivation as hybrid events. A state
crossing downward through the 10-CFU/mL LOD is set to zero and is hidden from
autonomous growth and outgoing-transfer terms. Incoming transfer from an
active paired compartment may accumulate until an upward crossing reaches the
LOD, at which point that state is reactivated. Every returned burden is either
zero or at least 10 CFU/mL. Solver failure, warnings, non-finite output, or
negative states raise an explicit integration error.

\section{Script execution (\texttt{if \_\_name\_\_ == "\_\_main\_\_"})}

Not a function -- this is top-level code that runs only when
\texttt{model\_Bacteremia.py} is executed directly (not when it is
imported). It builds one \texttt{PharmacokineticModel} and one
\texttt{ImmuneResponse}, defines a single baseline \texttt{params}
dictionary and initial condition $Y_0 = [0,\,0,\,100,\,100]$, integrates
\texttt{dual\_reservoir\_model} once over a 1944-hour (81-day) course with
the shared \texttt{integrate\_model} routine, and prints two sanity-check
values: the net
pre-treatment blood growth rate (asserted positive, to catch an
accidentally unrepresentative parameter set in which the untreated host
would clear infection on its own) and the peak sensitive-strain blood count
reached before vancomycin starts.

\section{Baseline parameter values}

Table~\ref{tab:params} lists every key in the \texttt{params} dictionary
consumed by \texttt{dual\_reservoir\_model}, with the baseline value used
in this project's scripts (values are consistent across scripts except
where noted).

\begin{table}[H]
\centering
\small
\renewcommand{\arraystretch}{1.15}
\caption{\texttt{params} dictionary keys used by \texttt{dual\_reservoir\_model}.}
\vspace{0.5cm}
\begin{tabularx}{\textwidth}{l l X}
\hline
\textbf{Key} & \textbf{Baseline value} & \textbf{Meaning} \\
\hline
\texttt{rho\_S} & 1.47 h$^{-1}$ & blood sensitive-strain growth rate (Vera-Yunca et al., in vitro) \\
\texttt{rho\_R} & 1.323 h$^{-1}$ & blood resistant-strain growth rate (10\% fitness cost relative to \texttt{rho\_S}) \\
\texttt{rho\_res\_S} & 0.845 h$^{-1}$ & reservoir sensitive-strain growth rate (Vera-Yunca et al., in vivo) \\
\texttt{rho\_res\_R} & 0.7605 h$^{-1}$ & reservoir resistant-strain growth rate (10\% fitness cost relative to \texttt{rho\_res\_S}) \\
\texttt{Emax\_v} & 1.7 h$^{-1}$ & vancomycin maximum kill rate \\
\texttt{EC50\_V} & 0.25 mg/L & vancomycin half-maximal effective concentration \\
\texttt{Emax\_l} & 1.97 h$^{-1}$ & linezolid maximum kill rate (Vera-Yunca et al., in vitro) \\
\texttt{EC50\_L} & 0.56 mg/L & linezolid half-maximal effective concentration (Vera-Yunca et al., in vitro) \\
\texttt{B\_max\_blood} & $6.0\times10^3$ CFU/mL & blood carrying capacity (shared by $S_b+R_b$) \\
\texttt{B\_max\_reservoir} & $1.0\times10^7$ CFU/mL & reservoir carrying capacity (shared by $S_{res}+R_{res}$) \\
\texttt{van\_res\_fraction} & 0.15 & vancomycin reservoir tissue-penetration fraction \\
\texttt{lzd\_res\_fraction} & 0.50 & linezolid bone-penetration fraction, midpoint of the reported 0.40--0.60 range \\
\texttt{lzd\_res\_efficacy} & 0.45 & reservoir linezolid efficacy factor $\alpha$ ($E_{max,L}^{res} = \alpha\,E_{max,L}$); calibrated \\
\texttt{k\_r\_b} & $1.0\times10^{-7}$ h$^{-1}$ & reservoir $\to$ blood exchange rate \\
\hline
\end{tabularx}
\label{tab:params}
\end{table}

\section{Parameters varied with a coefficient of variation}

\texttt{model\_Bacteremia.py} itself defines only fixed baseline values --
it contains no random sampling. Three quantities are varied in the current
analysis by separate Monte Carlo scripts that import
\texttt{dual\_reservoir\_model} and call it $n=1000$ times per script, each
call using one freshly sampled value. Log-normal distributions are
mean-corrected so their \emph{mean} (not their median) equals the baseline
value shown in Table~\ref{tab:params}. The $E_{max,L}$ analysis uses the same
mean-corrected log-normal parameterization, so its baseline value is likewise
the distribution's arithmetic mean. Every other parameter is held fixed at
baseline. Table~\ref{tab:cv} lists each sampled quantity's mean, CV, and
source script.

\begin{table}[H]
\centering
\small
\renewcommand{\arraystretch}{1.15}
\caption{Monte Carlo sampling distributions used elsewhere in this project
($n=1000$ draws each, mean-corrected).}
\vspace{0.5cm}
\begin{tabularx}{\textwidth}{l l l X}
\hline
\textbf{Quantity} & \textbf{Mean} & \textbf{CV} & \textbf{Source script} \\
\hline
$EC_{50,L}$ & 0.56 mg/L & 0.60 & \texttt{MC\_ec50\_lzd.py} \\
$E_{max,L}$ & 1.97 h$^{-1}$ & 0.080 & \texttt{MC\_emax\_lzd.py}; mean-corrected log-normal \\
$k_{immune}$ & 0.18 h$^{-1}$ & 0.80 & \texttt{MC\_immune\_response.py} \\
\hline
\end{tabularx}
\label{tab:cv}
\end{table}

The pharmacodynamic and immune CVs in Table~\ref{tab:cv} are
outcome-calibrated model inputs, not independently estimated measures of
population variability.

The sampled means for $EC_{50,L}$, $E_{max,L}$, and $k_{immune}$ equal their
fixed baseline values; in particular, \texttt{ImmuneResponse}'s default is
$0.18$ h$^{-1}$. Initial-condition sampling is excluded from the current
analysis. Within each one-factor Monte Carlo analysis, every non-target
parameter remains fixed at baseline. Separate deterministic sensitivity
scripts may deliberately sweep a named parameter, but those sweeps do not
assign it a population CV. Inter-compartmental transfer rates are held fixed.

\end{document}
